\section{模型总结与迁移}
\ifteachernotes
\block{三个条件，三条路线}
\begin{center}\small
\renewcommand{\arraystretch}{1.3}
\begin{tabularx}{\linewidth}{@{}p{.26\linewidth}p{.22\linewidth}X@{}}
\toprule
已知／目标 & 辅助线 & 真正起作用的关系\\
\midrule
两等腰直角，证明等面积 & 以等长边作底，分别作高 & 等底＋直角三角形全等得到等高\\
\addlinespace
已知中点，证明垂直 & 倍长 \(BE\) & 八字形全等搬运边，再比较大三角形\\
\addlinespace
已知垂直，证明中点 & 从 \(C\)、\(D'\) 向 \(BE\) 作双垂 & 两次一线三等角，再接八字形全等\\
\bottomrule
\end{tabularx}
\end{center}

在本章构型下，可以把题 2、题 3 合起来记：
\[
E\in CD':\qquad CE=ED'\ \Longleftrightarrow\ BE\perp C'D.
\]
题 2 还告诉我们：当 \(E\) 是 \(CD'\) 的中点时，\textbf{\(BE\) 既垂直于 \(C'D\)，又等于 \(C'D\) 的一半}。

\block{再看一次等面积}
题 3 的两组全等给出 \(CM=D'N=BF\)。以 \(BE\) 为底，\(\triangle CBE\) 和 \(\triangle D'BE\) 的高都是 \(BF\)，所以
\begin{align*}
S_{\triangle CBD'}
&=S_{\triangle CBE}+S_{\triangle D'BE}\\
&=BE\cdot BF.
\end{align*}
而题 2 已经得到 \(2BE=C'D\)，故
\[
S_{\triangle CBD'}=\frac12 C'D\cdot BF=S_{\triangle C'BD}.
\]
这给出等面积的另一种看法：倍长中线得到的长度关系，和双垂得到的高度关系，最终在面积公式里相遇。

\block{易错提醒}
\begin{enumerate}
\item 两个等腰直角三角形的方向要按图取定。只记两组等长和两个直角，不看方位，不能直接套用本章结论。
\item 高是到\textbf{直线}的距离，垂足可能在边的延长线上；题 1 的 \(N\) 就在 \(BC'\) 的延长线上。
\item 题 3 中的中点是结论，不能预先使用 \(CE=ED'\) 来证明全等。
\item “两块面积始终相等”不表示“每块面积是一个固定值”；构型变化时，两块面积可以一起变化。
\item 三题分别作图。题 1、题 3 的 \(M\)、\(N\) 含义不同；题 2 的 \(G\) 是倍长点，\(H\) 是交点。
\end{enumerate}
\else
\block{学生自主总结}
在每道题旁记录：已知条件、第一条辅助线、第一次全等搬运了什么、第二次全等得到什么。再尝试写出中点、垂直与长度之间的关系。
\vspace{8\baselineskip}
\fi
